% ============================================================================
% AERIS Paper - Section 3: System Model and Preliminary Analysis
% Target Journal: Applied Intelligence (APIN) - Springer
% ============================================================================

\section{System Model and Preliminary Analysis}
\label{sec:system}

This section presents the system model underlying AERIS, including network
assumptions, energy consumption model, and channel characteristics. We then
describe preliminary experiments that informed key design decisions.

\subsection{Network Model}

We consider a wireless sensor network consisting of $N$ sensor nodes
deployed in a two-dimensional area $A = W \times H$ square meters. The
network topology is characterized by the following assumptions:

\begin{enumerate}
    \item \textbf{Static deployment}: Nodes are stationary after initial
    deployment, a common assumption for environmental monitoring and
    industrial sensing applications.

    \item \textbf{Homogeneous nodes}: All sensor nodes have identical
    hardware capabilities (transmission power, sensing range, initial
    energy). We use CC2420-based TelosB motes as the reference platform.

    \item \textbf{Single base station}: A base station (BS) with unlimited
    energy is located at coordinates $(x_{BS}, y_{BS})$, typically outside
    the sensing area to simulate realistic deployment scenarios.

    \item \textbf{Data generation}: Each node generates one data packet
    per round, representing periodic sensor readings. Packet size is
    fixed at $L = 4000$ bits (500 bytes).

    \item \textbf{Bidirectional links}: If node $i$ can communicate with
    node $j$, then $j$ can also communicate with $i$ (symmetric links).
\end{enumerate}

The network is modeled as a graph $G = (V, E)$, where $V$ is the set of
sensor nodes and $E$ is the set of communication links. An edge $(i, j)
\in E$ exists if the Euclidean distance $d_{ij} \leq R_{max}$, where
$R_{max}$ is the maximum communication range.

\subsection{Energy Consumption Model}

We adopt an energy model based on the CC2420 radio transceiver, consistent
with IEEE 802.15.4 specifications \cite{cc2420datasheet}. The model accounts
for transmission, reception, and idle listening energy consumption.

\subsubsection{Transmission Energy}

The energy consumed to transmit $L$ bits over distance $d$ is:
\begin{equation}
E_{tx}(L, d) = E_{elec} \cdot L + E_{amp}(d) \cdot L
\label{eq:tx_energy}
\end{equation}

where $E_{elec} = 50$ nJ/bit represents the transceiver electronics energy,
and $E_{amp}(d)$ is the amplifier energy that depends on the propagation
model:
\begin{equation}
E_{amp}(d) = \begin{cases}
    E_{fs} \cdot d^2 & \text{if } d < d_0 \\
    E_{mp} \cdot d^4 & \text{if } d \geq d_0
\end{cases}
\label{eq:amp_energy}
\end{equation}

Here $E_{fs} = 10$ pJ/bit/m$^2$ is the free-space model coefficient,
$E_{mp} = 0.0013$ pJ/bit/m$^4$ is the multi-path model coefficient, and
$d_0 = \sqrt{E_{fs}/E_{mp}} \approx 87$ m is the crossover distance.

\subsubsection{Reception Energy}

The energy consumed to receive $L$ bits is:
\begin{equation}
E_{rx}(L) = E_{elec} \cdot L
\label{eq:rx_energy}
\end{equation}

\subsubsection{Data Aggregation Energy}

Cluster heads perform data aggregation (fusion) before transmission:
\begin{equation}
E_{agg}(L) = E_{DA} \cdot L
\label{eq:agg_energy}
\end{equation}

where $E_{DA} = 5$ nJ/bit is the aggregation energy per bit.

Table~\ref{tab:energy_params} summarizes the energy model parameters.

\begin{table}[htbp]
\centering
\caption{Energy Model Parameters (CC2420 TelosB)}
\label{tab:energy_params}
\begin{tabular}{@{}llc@{}}
\toprule
Parameter & Description & Value \\
\midrule
$E_{elec}$ & Electronics energy & 50 nJ/bit \\
$E_{fs}$ & Free-space amplifier & 10 pJ/bit/m$^2$ \\
$E_{mp}$ & Multi-path amplifier & 0.0013 pJ/bit/m$^4$ \\
$E_{DA}$ & Data aggregation & 5 nJ/bit \\
$d_0$ & Crossover distance & 87 m \\
$L$ & Packet size & 4000 bits \\
\bottomrule
\end{tabular}
\end{table}

\subsection{Channel Model}

We employ a log-normal shadowing model consistent with IEEE 802.15.4
propagation characteristics in indoor/industrial environments.

\subsubsection{Path Loss Model}

The received signal power at distance $d$ from a transmitter with power
$P_{tx}$ is:
\begin{equation}
P_{rx}(d) = P_{tx} - PL(d_0) - 10 \cdot \eta \cdot \log_{10}\left(\frac{d}{d_0}\right) + X_\sigma
\label{eq:path_loss}
\end{equation}

where $PL(d_0)$ is the path loss at reference distance $d_0 = 1$ m,
$\eta$ is the path loss exponent (environment-dependent), and $X_\sigma$
is a zero-mean Gaussian random variable with standard deviation $\sigma$
representing shadowing effects.

\subsubsection{Environment-Dependent Parameters}

Based on empirical measurements from the Intel Berkeley Research Lab
dataset \cite{madden2004intel} and standard IEEE 802.15.4 characterizations,
we define environment-specific channel parameters:

\begin{table}[htbp]
\centering
\caption{Channel Model Parameters by Environment Type}
\label{tab:channel_params}
\begin{tabular}{@{}lccc@{}}
\toprule
Environment & Path Loss $\eta$ & Shadowing $\sigma$ (dB) & Typical PDR \\
\midrule
Indoor Office & 2.0--2.5 & 3--4 & 95--99\% \\
Industrial & 2.5--3.5 & 4--8 & 85--95\% \\
Outdoor Open & 2.0--2.2 & 2--3 & 97--99\% \\
Urban Dense & 3.0--4.0 & 6--10 & 75--90\% \\
\bottomrule
\end{tabular}
\end{table}

\subsubsection{Packet Delivery Probability}

The probability of successful packet delivery over a link with distance
$d$ is modeled as:
\begin{equation}
P_{success}(d) = \left(1 - P_{BER}(SNR(d))\right)^L
\label{eq:pdr}
\end{equation}

where $P_{BER}$ is the bit error rate as a function of signal-to-noise
ratio, and $L$ is the packet length in bits.

\subsection{Preliminary Experiments and Design Rationale}

Before presenting the AERIS protocol, we describe three preliminary
experiments conducted on the Intel Berkeley Research Lab dataset that
informed key design decisions. This dataset comprises 2.22 million sensor
readings from 54 nodes collected over 36 days, including temperature,
humidity, light, and voltage measurements alongside link quality metrics.

\subsubsection{E1: Environment-Link Quality Correlation}

\textbf{Objective}: Quantify the relationship between environmental
parameters and wireless link quality.

\textbf{Method}: We computed Pearson correlation coefficients between
environmental variables (temperature, humidity) and link quality
indicators (RSSI, packet reception rate) across all node pairs.

\textbf{Results}:
\begin{itemize}
    \item Temperature--RSSI correlation: $r = -0.292$ ($p < 0.001$)
    \item Humidity--packet loss correlation: $r = 0.187$ ($p < 0.001$)
    \item Combined environment features predict link quality with
    classification accuracy of 89.2\%
\end{itemize}

\textbf{Design Decision $\rightarrow$ Environment-Aware Gateway Selection}:
The significant correlation between environmental conditions and link
quality motivates AERIS's environment-aware scoring function. By
incorporating temperature and humidity readings into gateway selection,
AERIS prioritizes nodes with stable environmental conditions, improving
aggregate link reliability.

\subsubsection{E2: Link Reliability Predictability}

\textbf{Objective}: Assess whether link reliability can be predicted
from observable node and environment features.

\textbf{Method}: We trained a random forest classifier using features
including node distance, residual energy, temperature, humidity, and
historical packet delivery rate to predict link success/failure.

\textbf{Results}:
\begin{itemize}
    \item Prediction AUC: 0.924 (10-fold cross-validation)
    \item Most important features: distance (42\%), temperature (18\%),
    historical PDR (15\%), humidity (12\%), energy (8\%), other (5\%)
\end{itemize}

\textbf{Design Decision $\rightarrow$ Hierarchical Backbone Routing}:
The high predictability of link reliability (AUC = 0.924) suggests that
routing decisions can be made deterministically based on observable
features, without requiring online learning. AERIS exploits this by
pre-selecting high-reliability links for the skeleton backbone, reserving
adaptive mechanisms (ARQ, cooperative relay) for edge cases.

\subsubsection{E3: Load Imbalance Impact}

\textbf{Objective}: Quantify the impact of traffic load imbalance on
network performance.

\textbf{Method}: We simulated various load distributions across cluster
heads and measured the correlation between load variance and end-to-end
packet delivery ratio.

\textbf{Results}:
\begin{itemize}
    \item Load variance--PDR correlation: $r = -0.749$ ($p < 0.001$)
    \item Networks with Gini coefficient $> 0.4$ showed 15--25\% PDR
    degradation
    \item Hotspot cluster heads depleted energy 3--5$\times$ faster
\end{itemize}

\textbf{Design Decision $\rightarrow$ Gateway Load Limiting}:
The strong negative correlation between load imbalance and PDR motivates
AERIS's gateway load limiting mechanism. By constraining the maximum
traffic each gateway can handle (\texttt{gateway\_load\_limit} parameter),
AERIS prevents hotspot formation and ensures balanced energy consumption
across cluster heads.

\subsection{Design Principles Summary}

Based on the preliminary analysis, AERIS is designed around three
core principles:

\begin{enumerate}
    \item \textbf{Environment-awareness}: Incorporate measurable
    environmental factors into routing decisions to improve link
    selection reliability (motivated by E1).

    \item \textbf{Deterministic hierarchy}: Use predictable, high-quality
    links for backbone routing rather than relying on computationally
    expensive online learning (motivated by E2).

    \item \textbf{Load balancing}: Actively limit gateway load to prevent
    hotspots and ensure uniform energy depletion (motivated by E3).
\end{enumerate}

These principles collectively enable AERIS to achieve $O(\log n)$ latency
through hierarchical routing while maintaining high reliability through
environment-aware link selection and load-balanced gateway coordination.
